% Full proofs of Theorem 1 (identification) and Theorem 2 (decomposition
% consistency) for the representation-derived-treatment framework.
% STATUS: self-verified draft. Algebraic cores checked exactly in
% data/scripts/theorem_symbolic_verify.py. The two previously-flagged steps are now
% verified: Thm 1 Step 6 (necessity) reduces to the exact bias beta*b*d/(d^2+1);
% Thm 2 Step 4 second-order coincidence confirmed numerically (gap ~ t^2) with the
% justification corrected from a Taylor argument to Neyman orthogonality of the
% Robinson moment. Compiles into the paper with zero undefined refs/cites. NOT yet
% \input into the committed manuscript, pending the author's final read.

\subsection{Identification of a representation-derived treatment}
\label{sec:thm-id}

Let $\ell$ denote location, $x$ the structured attributes, and $W$ the text
embedding, and let the scalar treatment $T$ be a fixed linear summary of $W$, the
oriented leading pooled principal component. Write $L=L(\ell)$ for a basis of
functions of location used as controls. We take the outcome to follow the
structural model
\begin{equation}
Y=\theta\,T+\beta\,U+f(x)+\varepsilon,
\label{eq:struct}
\end{equation}
in which $U$ is an unobserved, spatially patterned driver of price.

\begin{assumption}[Spatial confounder]\label{as:id-U}
The unobserved confounder is location-measurable, $U=h(\ell)$ for a measurable $h$.
\end{assumption}
\begin{assumption}[Basis richness]\label{as:id-basis}
$h$ lies in the closed linear span of $L$, so that $\mathbb{E}[U\mid L,x]=U$ almost
surely. Concretely $L$ includes the discrete geographic identifiers (borough,
neighborhood, school, and street indicators) through which the encoder leaks
location; a coordinate-polynomial basis alone does not satisfy this when $h$ varies
discontinuously across geographic units.
\end{assumption}
\begin{assumption}[Residual overlap]\label{as:id-overlap}
$\mathbb{E}\big[(T-\mathbb{E}[T\mid L,x])^2\big]>0$.
\end{assumption}
\begin{assumption}[Exogeneity]\label{as:id-exog}
$\mathbb{E}[\varepsilon\mid T,\ell,x]=0$.
\end{assumption}

\begin{theorem}[Identification]\label{thm:id}
Under Assumptions~\ref{as:id-U}--\ref{as:id-exog}, the treatment coefficient in
\eqref{eq:struct} equals the Robinson functional
\[
\theta=\frac{\mathbb{E}\!\left[\tilde Y\,\tilde T\right]}{\mathbb{E}\!\left[\tilde T^{2}\right]},
\qquad
\tilde Y=Y-\mathbb{E}[Y\mid L,x],\quad \tilde T=T-\mathbb{E}[T\mid L,x].
\]
The concept-erasure diagnostic consistently estimates $\mathbb{E}[\tilde T^{2}]$, and
$\theta$ is point-identified if and only if $\mathbb{E}[\tilde T^{2}]>0$; as
$R^2_{T\sim(L,x)}\to1$ the denominator vanishes and identification fails.
\end{theorem}

\begin{proof}
\emph{Step 1 (the confounding is absorbed).} By
Assumptions~\ref{as:id-U}--\ref{as:id-basis}, $U=h(\ell)$ with $h$ in the span of
$L$, so $\beta U+f(x)=g(L,x)$ for the fixed $(L,x)$-measurable function
$g:=\beta h+f$. Model \eqref{eq:struct} becomes $Y=\theta T+g(L,x)+\varepsilon$.

\emph{Step 2 (exogeneity survives coarsening).} Since $(L,x)$ is
$\sigma(\ell,x)$-measurable, iterated expectations and
Assumption~\ref{as:id-exog} give
$\mathbb{E}[\varepsilon\mid T,L,x]=\mathbb{E}\!\left[\mathbb{E}[\varepsilon\mid
T,\ell,x]\mid T,L,x\right]=0$.

\emph{Step 3 (partialling out).} Taking $\mathbb{E}[\,\cdot\mid L,x]$ of the model,
$\mathbb{E}[Y\mid L,x]=\theta\,\mathbb{E}[T\mid L,x]+g(L,x)$, and subtracting yields
$\tilde Y=\theta\,\tilde T+\varepsilon$.

\emph{Step 4 (the moment).} Multiplying by $\tilde T$ and taking expectations,
$\mathbb{E}[\tilde Y\tilde T]=\theta\,\mathbb{E}[\tilde T^{2}]+\mathbb{E}[\varepsilon
\tilde T]$. The cross term vanishes: $\mathbb{E}[\varepsilon\tilde
T]=\mathbb{E}[\varepsilon T]-\mathbb{E}[\varepsilon\,\mathbb{E}[T\mid L,x]]=0$, the
first term by $\mathbb{E}[\varepsilon T]=\mathbb{E}[\mathbb{E}[\varepsilon\mid
T,L,x]\,T]=0$ (Step 2) and the second by
$\mathbb{E}[\varepsilon\,\mathbb{E}[T\mid L,x]]=\mathbb{E}[\mathbb{E}[\varepsilon\mid
L,x]\,\mathbb{E}[T\mid L,x]]=0$. Dividing by $\mathbb{E}[\tilde T^{2}]$, positive by
Assumption~\ref{as:id-overlap}, gives the stated functional.

\emph{Step 5 (erasure diagnostic and the boundary).} $\mathbb{E}[\tilde
T^{2}]=\operatorname{Var}(T-\mathbb{E}[T\mid L,x])$ is the treatment's variance
orthogonal to $(L,x)$, which the held-out residual of the concept-erasure diagnostic
estimates consistently. Writing $R^2_{T\sim(L,x)}=1-\mathbb{E}[\tilde
T^{2}]/\operatorname{Var}(T)$, as this share tends to one the Robinson denominator
tends to zero and $\theta$ ceases to be point-identified. New York, whose leading
text direction loads on borough and neighborhood names, is the empirical instance of
this boundary.

\emph{Step 6 (necessity of Assumption~\ref{as:id-basis}).} If instead
$h\notin\operatorname{span}L$, let $r=U-\mathbb{E}[U\mid L]$ be the unabsorbed
remainder, which is $(L,x)$-orthogonal yet correlated with $\tilde T$ because $T$
shares the discrete location channel. Then $\tilde Y=\theta\tilde T+\beta\tilde
r+\varepsilon$ and the Robinson functional returns $\theta+\beta\,\mathbb{E}[\tilde
r\tilde T]/\mathbb{E}[\tilde T^{2}]\neq\theta$. Modeling a smooth channel
$\ell_s\in\operatorname{span}L$ and a discrete channel $\ell_d\notin
\operatorname{span}L$ with $U=a\ell_s+b\ell_d$ and $T=c\ell_s+d\ell_d+\eta$ gives the
exact bias $\beta\,b\,d/(d^{2}+1)$, zero if and only if the discrete channel is
absent from either $U$ or $T$. This is why the basis must span the discrete
geographic identifiers rather than coordinates alone.
\end{proof}

\subsection{The decomposition estimates the location-confounding bias}
\label{sec:thm-decomp}

Let $W_0=(x,c)$ be the base control set and $W_1=W_0\cup L$ the augmented set that
adds the location basis. Write $\hat\theta_{\mathrm{base}}$ and
$\hat\theta_{\mathrm{geo}}$ for the cross-fitted partially-linear DML coefficients on
$T$ under $W_0$ and $W_1$, and $\hat\Delta=\hat\theta_{\mathrm{base}}-\hat\theta_{
\mathrm{geo}}$.

\begin{assumption}[Cross-fitting rates]\label{as:dml-rates}
The nuisance estimators under both control sets satisfy the conditions of
\citet{chernozhukov2018double}, with product errors $o_p(n^{-1/2})$ and finite
second moments.
\end{assumption}
\begin{assumption}[Positivity]\label{as:dml-overlap}
$\mathbb{E}[\tilde T_0^{2}]$ and $\mathbb{E}[\tilde T_1^{2}]$ are bounded away from
zero, where $\tilde T_j$ residualizes $T$ on $W_j$.
\end{assumption}

\begin{theorem}[Decomposition consistency]\label{thm:decomp}
Under Assumptions~\ref{as:dml-rates}--\ref{as:dml-overlap}:
\begin{enumerate}
\item[\textup{(i)}] $\hat\Delta$ is asymptotically linear with influence function
$\psi_{\mathrm{base}}-\psi_{\mathrm{geo}}$, so $\hat\Delta\xrightarrow{p}\Delta:=
\theta_{\mathrm{base}}-\theta_{\mathrm{geo}}$ and
$\sqrt n(\hat\Delta-\Delta)\xrightarrow{d}\mathcal N\!\big(0,\
\mathbb{E}[(\psi_{\mathrm{base}}-\psi_{\mathrm{geo}})^{2}]\big)$.
\item[\textup{(ii)}] When the first-stage projections of $Y$ and $T$ on the controls
are linear, $\Delta$ equals the \citet{gelbach2016covariates} order-invariant
decomposition term for the location block, $\Delta=\gamma^{\top}\pi$, with $\gamma$
the location coefficient in the full linear projection and
$\pi_k=\mathbb{E}[\tilde T_0\tilde L_{0,k}]/\mathbb{E}[\tilde T_0^{2}]$.
\item[\textup{(iii)}] In general $\Delta$ equals the omitted-variable-bias functional
of \citet{chernozhukov2022long} for the location block, coinciding with the linear
expression of \textup{(ii)} to second order in the first-stage approximation error by
Neyman orthogonality of the Robinson moment.
\end{enumerate}
\end{theorem}

\begin{proof}
\emph{Step 1 (asymptotic linearity of each coefficient).} Under
Assumptions~\ref{as:dml-rates}--\ref{as:dml-overlap} and
\citet[Theorems~3.1--3.2]{chernozhukov2018double}, each cross-fitted DML coefficient
is asymptotically linear, $\sqrt n(\hat\theta_j-\theta_j)=n^{-1/2}\sum_i\psi_{j,i}+
o_p(1)$ for $j\in\{\mathrm{base},\mathrm{geo}\}$, with the Robinson influence function
$\psi_{j,i}=\mathbb{E}[\tilde T_j^{2}]^{-1}\,\tilde T_{j,i}\,(\tilde Y_{j,i}-\theta_j
\tilde T_{j,i})$.

\emph{Step 2 (the difference).} A difference of asymptotically linear estimators is
asymptotically linear with influence function the difference, so $\sqrt
n(\hat\Delta-\Delta)=n^{-1/2}\sum_i(\psi_{\mathrm{base},i}-\psi_{\mathrm{geo},i})+
o_p(1)$; consistency and the central limit theorem follow, giving (i).

\emph{Step 3 (linear case equals Gelbach).} Under linear first stages, nested-projection
algebra gives the omitted-variable-bias identity $\theta_{\mathrm{base}}=
\theta_{\mathrm{geo}}+\gamma^{\top}\pi$, with $\gamma$ the location coefficient in the
long projection and $\pi_k=\mathbb{E}[\tilde T_0\tilde
L_{0,k}]/\mathbb{E}[\tilde T_0^{2}]$ the coefficient of $L_k$ on $T$ after partialling
$W_0$. This is exactly the \citet{gelbach2016covariates} share for a single covariate
block, and his order-invariance result shows the assigned share is independent of the
sequence in which blocks enter. A symbolic computation confirms
$\theta_{\mathrm{base}}-\theta_{\mathrm{geo}}=(c\gamma_s+d\gamma_d)/(c^{2}+d^{2}+1)$
equals $\gamma^{\top}\pi$ exactly in the two-channel model, giving (ii).

\emph{Step 4 (general case equals the nonparametric OVB functional).} Without
linearity, $\Delta$ is the difference of two Robinson estimands under nested control
sets, the omitted-variable bias from dropping $L$, whose representation for a
partially-linear DML functional is that of \citet{chernozhukov2022long}. The linear
realization of Step 3 recovers this to second order in the first-stage approximation
error, and not merely to first order, because the Robinson moment is Neyman-orthogonal
to the first-stage projections: since $\mathbb{E}[\tilde T\mid W]=0$ by construction,
the Gateaux derivative of the moment in the outcome and treatment nuisances vanishes
at the truth \citep{chernozhukov2018double}, so replacing the true conditional
expectations by their linear projections perturbs each estimand only at second order.
A numerical check confirms the gap scales as the square of the first-stage
nonlinearity. This gives (iii).
\end{proof}
